
\documentclass[10pt,reqno]{amsart}
\usepackage[margin=0.55in]{geometry}
\usepackage{amsmath,amssymb,bbm,booktabs,graphicx,hyperref}
\usepackage{hyperref}
\usepackage{booktabs}
\usepackage{subcaption} % Required in your preamble
\hypersetup{
    colorlinks=true,
    linkcolor=blue,
    filecolor=magenta,      
    urlcolor=cyan,
}

\newcommand{\Pbb}{\mathbb P}
\newcommand{\R}{\mathbb R}
\newcommand{\T}{\mathcal T}
\newcommand{\ind}{\mathbbm 1}
\newcommand{\eps}{\varepsilon}

\title{Milestone 4: Critical Evaluation and Revision of the FraudYelp Ising Model (\href{https://github.com/Sanchayan-Bhowal/BitcoinTransactionClassification}{code link})}
\author{Tian Bai \and Sanchayan Bhowal \and Nathan Nguyen}
\date{}

\begin{document}
\maketitle
\vspace{-1.2em}

\section{Model}
We encode the label as a spin variable $\sigma_i \in \{-1,+1\}$, with $-1$ illicit and $+1$ licit. Let $X_i \in \mathbb{R}^d$ be the node feature vector, and let $E_1, E_2, E_3$ be the three relation-specific edge sets.
$\texttt{net\_rur},\texttt{net\_rsr},\texttt{net\_rtr}$.  For relation $r$,
\[
    w_{ij}^{(r)}
    =
    \exp\left(-\frac{\|X_i-X_j\|_2^2}{2\tau_r^2}\right),
    \qquad
    D_i^{(r)}=\sum_k A_{ik}^{(r)}w_{ik}^{(r)}.
\]
The degree-normalized interaction weight is
\[
    J_{ij}^{(r)}
    =
    A_{ij}^{(r)}w_{ij}^{(r)}
    \left(\frac{1}{D_i^{(r)}}+\frac{1}{D_j^{(r)}}\right).
\]
The fitted multi-relation Ising model is
\[
\Pbb_{\beta,\theta}(\sigma\mid X,G_1,G_2,G_3)
\propto
\exp\left\{
    \sum_{r=1}^3\beta_r\sum_{i,j}J_{ij}^{(r)}\sigma_i\sigma_j
    +
    2\sum_i\sigma_i b_\theta(X_i)
\right\},
\]
where $b_\theta$ is either linear or an MLP. Since the partition function is
intractable, we fit the model by pseudo-likelihood and optimize it with AdamW.

We use the training labels $\sigma_{\T}$, to get the estimates $\hat{\beta}$ and $\hat{\theta}$. As part of Milestone 3 we estimated the fraud probability as:
\[
    p_i^{\mathrm{cond}}
    =
    \Pbb_{\widehat\beta,\widehat\theta}(\sigma_i=-1\mid \sigma_\T,X,G)
    =
    \operatorname{logit}^{-1}(-2\widehat h_i).
\]

\section{Criticism and Revision: MCMC Marginals}

The main criticism is that $p_i^{\mathrm{cond}}$ is a local conditional
probability, whereas the prediction target under the fitted dependent-label
model is the posterior marginal
\[
    p_i^{\mathrm{marg}}
    =
    \Pbb_{\widehat\beta,\widehat\theta}(\sigma_i=-1\mid X,G).
\]
This marginal averages over the unknown labels outside the training set. The
conditional score does not perform this averaging and therefore does not fully
use the joint Ising model.

Our revision is to estimate posterior marginals by Gibbs sampling. The Gibbs update uses
\[
    H_i(\sigma_{-i})
    =
    \sum_{r=1}^3
    \widehat\beta_r\sum_{j\ne i}J_{ij}^{(r)}\sigma_j
    +
    2b_{\widehat\theta}(X_i),
    \qquad
    \Pbb(\sigma_i=-1\mid\sigma_{-i},X,G)
    =
    \operatorname{logit}^{-1}(-2H_i).
\]
After burn-in and thinning, we estimate
\[
    \widehat p_i^{\mathrm{MCMC}}
    =
    \frac1M\sum_{m=1}^M\ind\{\sigma_i^{(m)}=-1\}.
\]
The final classifier thresholds $\widehat p_i^{\mathrm{MCMC}}$, with the
threshold chosen on validation F1.

\begin{table}[h]
\centering
\caption{F1 at the validation-selected best-F1 threshold, rounded from the
experiment plot.}
\vspace{-0.35em}
\begin{tabular}{lccc}
\toprule
Model & conditional score & MCMC marginal score & change \\
\midrule
\texttt{hm\_weighted\_nn}     & 0.710 & 0.738 & +0.027 \\
\texttt{gm\_weighted\_nn}     & 0.686 & 0.711 & +0.025 \\
\texttt{hm\_weighted\_linear} & 0.646 & 0.674 & +0.028 \\
\texttt{unweigh\_neural}      & 0.666 & 0.655 & -0.011 \\
\texttt{unweigh\_linear}      & 0.595 & 0.595 &  0.000 \\
\texttt{gm\_weighted\_linear} & 0.556 & 0.583 & +0.027 \\
\bottomrule
\end{tabular}
\end{table}

\vspace{-0.7em}
\section{Sensitivity Analysis}
\subsection{Edge-Type Ablation}
We evaluated sensitivity to the modeling assumption that all three relations
should be included. For each ablation we refit the model and evaluated on the
test split using the validation-F1 threshold.

\begin{table}[h]
\centering
\caption{Edge-type ablation results.}
\vspace{-0.35em}
\resizebox{\textwidth}{!}{
\begin{tabular}{llrrrr}
\toprule
Ablation & Edge types used & F1 & ROC-AUC & Avg. precision & TP/FN/FP/TN \\
\midrule
\texttt{all\_3\_edges}
& \texttt{rur,rsr,rtr}
& 0.695562 & 0.930746 & 0.766819 & 856/444/305/7587 \\
\texttt{no\_net\_rsr}
& \texttt{rur,rtr}
& 0.691794 & 0.927995 & 0.754971 & 881/419/366/7526 \\
\texttt{no\_net\_rtr}
& \texttt{rur,rsr}
& 0.688172 & 0.929968 & 0.762140 & 832/468/286/7606 \\
\texttt{only\_net\_rur}
& \texttt{rur}
& 0.682306 & 0.925771 & 0.749799 & 858/442/357/7535 \\
\texttt{only\_net\_rsr}
& \texttt{rsr}
& 0.619048 & 0.891669 & 0.672229 & 780/520/440/7452 \\
\texttt{no\_net\_rur}
& \texttt{rsr,rtr}
& 0.615948 & 0.892047 & 0.677019 & 757/543/401/7491 \\
\texttt{only\_net\_rtr}
& \texttt{rtr}
& 0.601473 & 0.887858 & 0.660800 & 735/565/409/7483 \\
\bottomrule
\end{tabular}}
\end{table}

The full model is best by F1 and average precision. However,
\texttt{only\_net\_rur} is close to the full model:
\[
    0.695562-0.682306=0.013256
\]
in F1. In contrast, removing \texttt{net\_rur} reduces F1 to $0.615948$.
Therefore the conclusions are sensitive to whether \texttt{net\_rur} is
included, but less sensitive to the inclusion of \texttt{net\_rsr} or
\texttt{net\_rtr} once \texttt{net\_rur} is present.

\subsection{Sensitivity of posterior marginal uncertainty}

As an additional sensitivity check, we compared the uncertainty of the MCMC
marginals across model variants. After estimating $(\widehat\beta,\widehat\theta)$,
we generated full MCMC samples
\[
    \sigma^{(1)},\ldots,\sigma^{(M)}\in\{-1,+1\}^N
\]
from the fitted Ising model. For each node, we estimated the marginal fraud
probability
\[
    \widehat p_i
    =
    \frac1M
    \sum_{m=1}^M
    \mathbbm 1\{\sigma_i^{(m)}=-1\}.
\]
We then examined the marginal Bernoulli variance
\[
    \widehat p_i(1-\widehat p_i).
\]
This quantity is largest near $\widehat p_i=1/2$ and close to zero when the
posterior marginal is near $0$ or $1$. Thus it measures posterior uncertainty in
the node-wise classification score.

\begin{figure}[h]
    \centering
    \includegraphics[width=0.95\textwidth]{marginal_variance_violin.png}
    \caption{Distribution of MCMC marginal variances
    $\widehat p_i(1-\widehat p_i)$ over non-training nodes.}
    \label{fig:marginal_variance}
\end{figure}

Figure~\ref{fig:marginal_variance} shows that the HM weighted model and the
unweighted model produce substantially smaller marginal variances than the
GM weighted model. Thus, under the fitted MCMC posterior, these models give more
concentrated marginal predictions. In particular, the HM weighted model is both
aligned with our specified degree-normalized weighted Ising likelihood and gives
lower posterior uncertainty than the GM weighted variants.

For the HM and unweighted models, the linear external-field versions have lower
marginal variance than the neural versions. This is expected because the linear
field is less flexible and produces smoother, less variable node-level logits.
The neural field can represent more heterogeneous covariate effects, which can
leave more nodes near the decision boundary and hence increase
$\widehat p_i(1-\widehat p_i)$. We do not interpret lower variance alone as
better; rather, we use it together with F1 and average precision. The main
conclusion is that the chosen HM weighted model gives competitive predictive
performance while avoiding the large posterior marginal uncertainty seen in the
GM weighted variants.

\section{Comparison}

The original model uses the conditional fraud score $p_i^{\mathrm{cond}}$.
The revised model uses the MCMC posterior marginal score
$\widehat p_i^{\mathrm{MCMC}}$. The revised score is better for the weighted
models by F1 at the validation-selected threshold, with gains of about
$0.025$--$0.028$ for the weighted neural and weighted linear variants. The
ablations show that the best edge specification remains the full three-relation
model, while \texttt{net\_rur} alone is a competitive simplification.

\end{document}
